\apendice{Orientações Gerais sobre Apêndices}
\label{ap:orientacoes-gerais-sobre-apendices}

Elemento opcional (Guia UECE 2026, 2.2.4.3). Apêndice é o texto ou documento elaborado pelo próprio autor que complementa a argumentação do trabalho sem sobrecarregar o corpo do texto -- por exemplo, o roteiro completo de uma entrevista, o questionário aplicado, o código-fonte extenso, ou dados brutos de uma coleta.

A diferença para o Anexo (ver \autoref{an:orientacoes-gerais-sobre-anexos}) é a autoria: o Apêndice é produzido pelo próprio autor do TCC; o Anexo é um documento de terceiros, incluído apenas como fundamentação, comprovação ou ilustração.

Regras de formatação que não são automáticas neste template e exigem atenção do autor:

\begin{alineas}
	\item cada apêndice é identificado por uma letra maiúscula consecutiva, travessão e o respectivo título -- por isso o comando \texttt{\textbackslash apendice\{título\}} deve receber um título diferente em cada arquivo, na ordem em que aparecem no documento.tex;
	\item quando as 26 letras do alfabeto se esgotarem, usam-se letras dobradas (AA, AB, AC...);
	\item no Sumário e no corpo do texto, a palavra APÊNDICE e a letra correspondente devem aparecer em letras maiúsculas e em negrito -- o macro já cuida disso automaticamente.
\end{alineas}

Para acrescentar um novo apêndice ao seu trabalho, crie um novo arquivo .tex nesta pasta e adicione a linha \texttt{\textbackslash input\{...\}} correspondente, na seção \texttt{\textbackslash imprimirapendices} do documento.tex.
